\documentclass[12pt]{article}
\usepackage[utf8]{inputenc}
\usepackage[T1]{fontenc}
\usepackage{hyperref}
\title{Yeyi Language Resources for K-12 Curriculum Development}
\author{Compiled by Claude Sonnet 5.5 (claude-sonnet-5-5) for Andr\'as Kornai}
\date{}
\begin{document}
\maketitle

\section{Introduction}
Yeyi (autonym Shiyeyi; also Yeei, Yei; ISO 639-3 \texttt{yey}, Glottocode yeyi1239) is a Bantu language spoken along the Okavango River in Botswana and Namibia, by many of the roughly 50,000 Yeyi people. It is known for the largest click inventory of any Bantu language. Its main dialect is Shirwanga. The Endangered Languages Project rates it severely endangered. It is being displaced by Tswana in Botswana and by Lozi in the Caprivi area of Namibia.

\section{Primary Educational Resources}
No primers, textbooks or school programmes were found in the search results.

\section{Linguistic and Reference Materials}
\begin{itemize}
\item Endangered Languages Project, Yeyi entry. Its sources conflict on orthography: one record reports an orthography agreed at a conference in April 1997, while others say there is no standardised or published orthography: \url{https://endangeredlanguages.com/lang/1886}
\item Word-list pages from the Endangered Languages Project (1998 word lists): \url{https://endangeredlanguages.com/node/86237}
\item Library of Congress record for Lukusa, S.T.M., \emph{Groundwork in Shiyeyi grammar}, including a Shiyeyi--English glossary (2002). The search did not show a direct URL for the work itself; the catalogue authority page is: \url{https://nli.org.il/en/authorities/987007530319505171}
\end{itemize}

\section{Religious and Community Resources}
No scripture or community organisation was found in the search results.

\section{Digital and Audio Resources}
No audio or video resources were found. The word lists above are the only digital lexical data seen.

\section{Curriculum Development Context}
Yeyi is not known to be taught in schools. The orthography is unsettled according to the sources found, so a first step would be to confirm the status of the 1997 orthography with Yeyi communities and linguists. Tswana, Lozi and English dominate schooling in the region. This survey found no Yeyi-language curriculum materials.
\end{document}